%% Türkçe makale taslağı
%% Derleme: pdflatex makale_tr.tex  (elsarticle.cls gerektirir)
%% Grafikler için yer tutucular [SEKIL X] biçiminde işaretlenmiştir.

\documentclass[preprint,12pt]{elsarticle}

\usepackage[turkish]{babel}
\usepackage{amsmath,amssymb}
\usepackage{graphicx}
\usepackage{booktabs}
\usepackage{siunitx}
\usepackage{xcolor}
\newcommand{\todo}[1]{{\color{red}\textbf{(YAPILACAK: #1)}}}

\journal{Physical Review Accelerators and Beams}

\begin{document}

\begin{frontmatter}

\title{Dondurulmuş-spin proton depolama halkasında sahte EDM'ye yol açan
baskın harmonik bileşenin tek-yörünge tahmini}

\author[inst1]{S.~Hac\i\"omero\u{g}lu\corref{cor1}}
\cortext[cor1]{Sorumlu yazar}
\ead{selcuk.haciomeroglu@gmail.com}

\address[inst1]{\todo{kurum bilgisi}}

\begin{abstract}
Proton elektrik dipol momenti (pEDM) dondurulmuş-spin deneyi, depolama
halkasındaki her kuadrupol mıknatısın düşey konumunun $\sim\SI{10}{\micro\meter}$
hassasiyetle bilinmesini gerektirmektedir; düzeltilmemiş hizalama hataları gerçek
bir EDM sinyalini taklit eden kapalı-yörünge bozunumuna (KYB) yol açar.
Işın konum monitörleri (BPM) $\sim\SI{100}{\micro\meter}$ mertebesinde statik
elektronik ofset taşır ve bu durum hizalama sinyalini gömebilir.
Yerleşik çözüm olan $K$-modülasyonu (k-mod) iki kafes yapılandırmasının farkını
alarak ofseti tam olarak giderir; ancak yalnızca birkaç kuadrupol modüle
edilebildiğinde fark matrisi $\Delta R$'nin rank'ı $N_\mathrm{mod}$'a
düşer ve koşullanma sayısı $\kappa(\Delta R)\sim10^6$'ya fırlar —
geri çatımı pratikte olanaksız kılar.
Bu çalışmada, 48 boyutlu hizalama vektörünü tam olarak geri çatmak yerine, yalnızca
baskın sahte-EDM harmonik bileşenini ($k=2$) tahmin etmeyi hedefleyen
bir yöntem sunulmaktadır. Tek bir kapalı-yörünge ölçümü, tepki matrisi Fourier
tabanında CLEAN algoritmasıyla çözümlenerek $k=1,2,3$ modları \SI{100}{\micro\meter}
BPM ofseti altında $\sim5\%$ göreli genlik doğruluğuyla kurtarılmaktadır; gradyan
modülasyonuna gerek yoktur. Bu başarı tepki matrisinin $R$ üç nicel özelliğine
dayanmaktadır: (i)~ayar frekansına yakınlıktan doğan rezonans kuvvetlenmesi
($\|RF_{k=2}\|$ kafes ortalamasının $\sim34$ katı), (ii)~bunun sonucu olarak
yörünge düzeyinde hizalama genliği hiyerarşisinin tersine dönmesi ve
(iii)~farklı harmoniklerin yörünge tepkilerinin neredeyse dikeyliği, bu sayede
büyük kirletici harmonikler $k=2$ tahminine sızmamaktadır.
Yöntem, kuadrupol kaynaklı hatalar kadar harici manyetik alan harmoniklerini de
aynı çerçevede kapsamakta; BBA sonrası $\sigma_b\lesssim\SI{20}{\micro\meter}$
koşulunda $k=2$ bileşeni $\sim\SI{1}{\nano\tesla}$ hassasiyetle ölçülebilmektedir.
\end{abstract}

\begin{keyword}
depolama halkası \sep elektrik dipol momenti \sep kapalı-yörünge bozunumu \sep
kuadrupol hizalaması \sep tepki matrisi \sep ışın konum monitörü \sep CLEAN algoritması
\end{keyword}

\end{frontmatter}

%==============================================================
\section{Giriş}
\label{sec:giris}

\subsection{Kuadrupol hizalama hatası ve sahte EDM mekanizması}

Proton elektrik dipol momenti (pEDM) arayışı dondurulmuş-spin modunda
çalıştırılır: düzlem içi spin presesyonu sıfırlanır ve gerçek bir EDM,
yavaş birikimli bir düşey spin dönmesine yol açar. Bu düzenekteki baskın
sistematik sınırlama, kuadrupol mıknatısların düşey kapalı yörünge bozunumudur
(KYB). $j$'inci kuadrupol $\Delta y_j$ kadar düşey yer değiştirirse ışını
yönlendirir; oluşan KYB, bir proton gerçek EDM'ye sahipmiş gibi spini döndüren
radyal alanlar örnekler.

Sayısal olarak, bu sahte sinyalin hedef hassasiyetin altında tutulması için
her kuadrupolün düşey konumunun $\sim\SI{10}{\micro\meter}$ düzeyinde bilinmesi
gerekir; bu hassasiyet mekanik ölçümün çok ötesinde olduğundan ışın-tabanlı
bir belirleme zorunludur.

Sahte EDM oranı ile hizalama hatası arasındaki bağıntı doğrusal ve zamandan
bağımsızdır. Dondurulmuş-spin geometrisinde, $j$'inci kuadrupolün yaptığı katkı
yerel kapalı-yörünge konumu $y_\mathrm{co}(s_j)$ ile orantılıdır:
\begin{equation}
  \frac{dS_y}{dt} \propto
  \sum_j K_j L_j \left[y_\mathrm{co}(s_j) - \Delta y_j\right],
  \label{eq:spinrate}
\end{equation}
burada $K_j$ ve $L_j$ $j$'inci kuadrupolün efektif gradyanı ve uzunluğu,
$y_\mathrm{co}(s_j)$ kapalı-yörünge konumu ve $\Delta y_j$ mekanik kaçıklığıdır.
$y_\mathrm{co}$ da $\Delta y_j$'ye bağlı olduğundan bu ifade $\Delta q$'da
gerçekte ikinci dereceye kadar gider; ancak küçük kaçıklıklar için lineer
yaklaşım yeterlidir ve $\Delta q$'nun Fourier bileşenleri cinsinden ayrışır.
$k$'ıncı Fourier bileşeni $\hat{a}_k$ ise Denklem~(\ref{eq:spinrate})'e
$\|\hat{a}_k\|$ ile orantılı katkı verir.
Önemli olan, bu katkının $k=2$
Fourier harmonik bileşeni için ($N_\mathrm{hücre}=24$ FODO hücreli halkada,
toplam $N_Q=2N_\mathrm{hücre}=48$ kuadrupol, düşey tune $Q_y\approx2{,}68$)
diğer tüm bileşenlerden bir mertebe fazla olmasıdır;
bunun nedeni $k=2$'nin rezonansa yakınlığıdır (bkz. Bölüm~\ref{sec:fourier}).
Şekil~\ref{fig:falseedm}'de görüldüğü gibi $dS_y/dt$ spektrumu $k=2$'de keskin bir
tepe sergilemekte, $k=1$'i $2{,}9$ kat, $k=3$'ü $2{,}2$ kat aşmaktadır.

Bu çalışmadaki tüm sayısal simülasyonlar, Omarov et al.~\cite{omarov}'deki
kapsamlı pEDM halka tasarımını esas almaktadır: $N_\mathrm{hücre}=24$ FODO
hücresi, yarıçap $R_0=\SI{95{,}49}{\meter}$, nominal kuadrupol gradyanı
$G=\SI{0{,}21}{\tesla\per\meter}$, $Q_y=2{,}68$.
Beta fonksiyonları, faz ilerlemesi ve tepki matrisi bu kafesin doğrusal optik
çözümünden sayısal olarak elde edilmiştir; pEDM tasarım çalışmasına özgü kafes
geometrisi kullanıldığından bu çalışmanın bulguları doğrudan o tasarıma
uygulanabilir niteliktedir.

\begin{figure}[ht]
\centering
\fbox{\parbox{0.75\textwidth}{\centering\vspace{2cm}
\textbf{[ŞEKİL 1]} \\
$k = 0$--$12$ için $|dS_y/dt|$ tarama sonuçları (Richardson hata barlarıyla). \\
Her veri noktası: tek harmonik kaçıklık $A=\SI{10}{\micro\meter}$,
stroboskopik spin takibi, $t_2=\SI{2}{\milli\second}$. \\
$k=2$'de belirgin tepe; kapalı-yörünge genliği (ikincil eksen) de $k=2$'de tepe.
\vspace{2cm}}}
\caption{Sahte-EDM oranı $|dS_y/dt|$ ile kapalı-yörünge genliğinin $k$'ya göre taraması.
Her ölçümde yalnızca ilgili $k$ bileşeni sıfırdan farklıdır ($A=\SI{10}{\micro\meter}$).
Oran $k=2$'de tepe yapar: $k=2$'de $dS_y/dt = 1{,}42\times10^{-9}$~rad/s;
$k\ge9$'da kapalı-yörünge genliği $<\SI{3}{\micro\meter}$'e inerek $\lesssim\!2\%$
düzeyinde kalmaktadır.}
\label{fig:falseedm}
\end{figure}

\subsection{Temel ölçüm denklemi ve BPM ofseti sorunu}

Ölçüm modeli doğrusaldır:
\begin{equation}
  \mathbf{y} = R\,\Delta q + \mathbf{b} + \boldsymbol{\eta},
  \label{eq:model}
\end{equation}
burada $\mathbf{y}\in\mathbb{R}^{48}$ BPM okumaları, $R$ $48\times48$ yörünge
tepki matrisi, $\Delta q$ bilinmeyen kuadrupol kaçıklıkları, $\mathbf{b}$
statik BPM elektronik ofsetleri ($\sim\SI{100}{\micro\meter}$, saatler-günler
boyunca sabit fakat belirlenemeyen) ve $\boldsymbol{\eta}\sim\SI{1}{\micro\meter}$
okuma gürültüsüdür. Güçlük tamamen $\mathbf{b}$'dedir: $R^{-1}\mathbf{y}$
üzerinden $\Delta q$'yu çözmek $R^{-1}\mathbf{b}$ kirletmesini taşır ve
$\kappa(R)\approx249$ ile bu kirletme milimetre mertebesine çıkar — hedefin
iki mertebe üstünde.

\subsection{Mevcut yaklaşımlar: LOCO ve $K$-modülasyonu}

\paragraph{LOCO (Lineer Optik Kapalı-Yörüngeden).}
LOCO, BPM'lerden ölçülen yörünge tepki matrisini kafes modeline iteratif
olarak uyduran standart bir algoritmadır~\cite{safranek1997}. Güçlü yönü
kuadrupol güçlenme hatalarını ($\delta K/K$) tek bir ölçüm setiyle
yakalayabilmesidir. Kısıtlamaları şunlardır:
\begin{itemize}
  \item \textbf{BPM ofseti:} $\mathbf{b}$ tüm yörünge ölçümlerine girer
    ve BPM kazanç/faz kalibrasyonu yapılmadan kaldırılamaz; bu durum
    hizalama doğruluğunu sınırlar.
  \item \textbf{BPM tilti:} Açısal hizalama hataları ölçüm matrisine
    sistematik sapma ekler ve tam olarak giderilmesi için ek çalışma
    gerektirir.
  \item \textbf{Ölçüm süresi:} Kapsamlı bir LOCO analizi birden fazla
    yörünge ölçümü ve veri seti gerektirir; gerçek zamanlı izleme için
    pratik kısıt oluşturur.
\end{itemize}

\paragraph{$K$-modülasyonu (k-mod).}
K-modülasyonda iki kafes yapılandırması ölçülür ve fark alınır:
\begin{equation}
  \Delta\mathbf{y} = \mathbf{y}_2 - \mathbf{y}_1
  = (R_2 - R_1)\,\Delta q + (\boldsymbol{\eta}_2 - \boldsymbol{\eta}_1)
  \equiv \Delta R\,\Delta q + \text{gürültü}.
\end{equation}
Bu işlem $\mathbf{b}$'yi tam olarak eler. Avantajları:
BPM ofseti artık sorun değildir; fark alma sonrası BPM tilti etkisi de
azalır. Kısıtlamalar ise şunlardır:

\begin{itemize}
  \item \textbf{Tüm kuadrupoller birlikte modüle edilmelidir.}
    Yalnızca birkaç kuadrupol modüle edilirse $\Delta R$'nin sütunları
    birbirine çok benzediğinden rank düşer ve koşullanma sayısı $\sim10^6$'ya
    fırlar; geri çatım ciddi biçimde bozulur. Buna karşın, tüm 48 kuadrupol
    birlikte modüle edildiğinde $\Delta R = \varepsilon R$ ve
    $\kappa(\Delta R) = \kappa(R) \approx 249$ olur — tam-halka k-modu
    kötü koşullanma sorununu yaratmaz. Bu, SVD analizi ile gösterilmiştir
    (Şekil~\ref{fig:svd}).
  \item \textbf{Operasyonel yük:} 48 kuadrupolün aynı anda hassas gradyan
    adımıyla modüle edilmesi gerçek bir işletme zorluğu oluşturur.
  \item \textbf{Süre kısıtı:} İki ayrı denge noktasında ölçüm almak ve
    veriyi işlemek zaman alır.
\end{itemize}

Sonuç olarak k-modülasyonu ile tam hizalama vektörünü kurtarmak ya \emph{tüm}
48 kuadrupolün eş zamanlı hassas modülasyonunu --- büyük operasyonel yük ---
ya da $\kappa(\Delta R)\sim10^6$ koşullanma cezasını göze almayı gerektirir.
Bu ikilemin somut karşılığı: modüle edilen kuadrupol sayısı $N_\mathrm{mod}<48$
olduğunda $\mathrm{rank}(\Delta R)\le N_\mathrm{mod}$; sistem eksik belirlenir
ve geri çatım bütünüyle çöker. Bu çalışmanın yazarları tam da bu sorunla
karşılaşmış, $N_\mathrm{mod}=4\ldots12$ için $\kappa(\Delta R)>10^6$
düzeyleri gözlemlemiştir.
\textbf{Bu rank kısıtı, mevcut çalışmanın birincil motivasyonudur}:
tam $\Delta q$'yu kurtarmak yerine yalnızca sahte-EDM açısından kritik
Fourier harmoniklerini ($k=1,2,3$) hedefleyerek rank sorununu yapısal
olarak aşıyoruz (Şekil~\ref{fig:kmod_rank}). Bunun yanı sıra aşağıda
açıklanan ortak körlük her iki yaklaşım için de geçerlidir.

\begin{figure}[ht]
\centering
\fbox{\parbox{0.78\textwidth}{\centering\vspace{2cm}
\textbf{[ŞEKİL 7]} \\
Sol: $k=2$ genlik tahmin hatası (\%) vs modüle edilen kuadrupol sayısı $N_\mathrm{mod}$,
kirletici modlar ($k=4\ldots10$, $100$--$300\,\mu$m) mevcut.
Yatay kesikli çizgi: tek yörünge CLEAN referansı ($\approx5\%$).
Sağ: $k=2$ sinyal yakalama oranı vs $N_\mathrm{mod}$.
Gerçek $N_\mathrm{mod}=3$ (dikey noktalı çizgi): hata $\approx4300\%$, yakalama $\approx7\%$.
\vspace{2cm}}}
\caption{Kısmi k-modülasyonunun rank kısıtı.
$N_\mathrm{mod}<48$ kuadrupol modüle edildiğinde $k=2$ sinyalinin yalnızca
$N_\mathrm{mod}/48$ fraksiyonu örneklenir ve kirleticiler mevcut olduğunda
tahmin hatası $>1000\%$'e fırlar.
Tam halka k-modu ($N_\mathrm{mod}=48$) hata $\approx1\%$ verir,
ancak 48 kuadrupolün eş zamanlı hassas modülasyonunu gerektirir.
Tek yörünge CLEAN yaklaşımı ise gradyan modülasyonuna gerek kalmadan
$\approx5\%$ hata düzeyinde kalır.}
\label{fig:kmod_rank}
\end{figure}

\subsection{Hem LOCO hem k-modülasyonunun ortak körlüğü}

LOCO ve k-modülasyonu ikisi de geometrik faz etkisini ve rezonans kaynaklı
sahte EDM sinyalini önemli ölçüde bastırabilir. Öte yandan her iki yöntem de
\textit{harici manyetik alan harmoniklerine karşı kördür}: halka boyunca yayılan
$B_x(\theta)=B_r\cos(k\theta)$ biçimindeki bir dış radyal alan, aynı Fourier
bileşeninde yörünge bozunumu yaratır ve kuadrupol hizalama hatasından
ayırt edilemez. Bu nedenle örgü boyunca genel $k=2$ modunun ölçülmesi ayrı bir
önem taşımaktadır: söz konusu ölçüm, mekanik hizalama hatalarını ve harici
alan harmoniklerini aynı çerçevede bir arada kapsar.

%==============================================================
\section{Fourier bileşen analizi ve tepki matrisi kazancı}
\label{sec:fourier}

\subsection{Hizalama hatası spektrumu ve $k=2$'nin baskınlığı}

Gerçek halkalarda kuadrupol hizalama hataları istatistiksel olarak dağılmıştır;
herhangi bir Fourier bileşeni diğerlerinden sistematik olarak büyük olmamaktadır.
Bununla birlikte \textit{ölçüm stratejisi} açısından $k=2$'ye odaklanmak
güçlü bir gerekçeye dayanmaktadır:

\begin{enumerate}
  \item Şekil~\ref{fig:falseedm}'de gösterildiği gibi sahte-EDM oranı $k=2$'de
    keskin biçimde tepe yapar; diğer bileşenlerin katkısı en az bir mertebe
    küçüktür. Sonraki en büyük bileşen $k=1$'dir ve yaklaşık dörtte biri
    kadardır.
  \item $k=2$ bileşeni tespit edildikten sonra halka boyunca dağılmış
    dipol düzeltici mıknatıslarıyla bastırılabilir; bu düzeltme planının
    tasarımı ayrı bir çalışma konusudur ve burada ele alınmamaktadır.
\end{enumerate}

\subsection{Tepki matrisi, eşleşmiş filtre kazancı ve SVD analizi}

Belirli bir $k$ bileşenindeki hizalama hatası, tepki matrisi üzerinden her
BPM'e Green fonksiyonu ağırlıklarıyla yansır:
\begin{equation}
  \mathbf{y}^{(k)} = R\,F_k\,\mathbf{a}_k,
\end{equation}
burada $F_k\in\mathbb{R}^{48\times2}$ FODO-antisimetrik Fourier tabanı ve
$\mathbf{a}_k=(a_{k,c},\,a_{k,s})^T$ cos/sin katsayılarıdır.
$k$'ıncı harmonik için eşleşmiş-filtre kazancı $\|M_k\|=\|RF_k[:,0]\|$
olarak tanımlanır.

\textbf{Kritik sayısal değer:} $\|M_{k=2}\|=167$, kafes ortalamasının
$\sim\!34$ katıdır. Bu, $k=2$'nin dikey ayar frekansına yakınlığından
($Q_y\approx2{,}68$) doğan rezonans kuvvetlenmesinin doğrudan
bir yansımasıdır.

Tepki matrisinin $\kappa(R)\approx249$ koşullanma sayısı bu yörünge kazancıyla
karıştırılmamalıdır. SVD analizi,
\begin{equation}
  R = U\,\Sigma\,V^T,
\end{equation}
$R$'nin en büyük iki tekil vektörünün tam olarak $k=2$ cos ve sin
modları olduğunu ortaya koymaktadır (Şekil~\ref{fig:svd}). Bu,
$k=2$'nin halkanın "doğal modu" olduğunun cebirsel kanıtıdır.

\begin{figure}[ht]
\centering
\fbox{\parbox{0.75\textwidth}{\centering\vspace{2cm}
\textbf{[ŞEKİL 2]} \\
Sol: $R$ tepki matrisinin SVD tekil değerleri ($\sigma_1,\ldots,\sigma_{48}$). \\
Sağ: En büyük iki sağ tekil vektör ($v_1$, $v_2$) —
her ikisi de tam $k=2$ cos/sin örüntüsüyle örtüşmektedir. \\
\vspace{2cm}}}
\caption{Tepki matrisi SVD analizi. En büyük iki tekil değer diğerlerinden
belirgin biçimde ayrışır; karşılık gelen tekil vektörler $k=2$ FODO-antisimetrik
modlarıdır. $\kappa(R)=\sigma_1/\sigma_{48}\approx249$ koşullanma sayısıdır;
bu rakam yörünge kazancı değildir.}
\label{fig:svd}
\end{figure}

\subsection{Üç genlik ölçeği ve BPM ofseti bastırımı}

$\SI{10}{\micro\meter}$'lik $k=2$ hizalama hatasını $\sigma_b=\SI{100}{\micro\meter}$
BPM ofseti arka planına karşı ölçmek için birbirinden farklı üç genlik ölçeği
devreye girmektedir (Şekil~\ref{fig:amplitudes}):

\begin{enumerate}
\item \textbf{Sinyal yörüngesi ($\SI{1669}{\micro\meter}$).}
  $A=\SI{10}{\micro\meter}$'lik $k=2$ hizalama hatası $R$ üzerinden
  rezonans kuvvetlenmesiyle $A\|M_{k=2}\|=10\times167=\SI{1669}{\micro\meter}$
  yörünge normuna dönüşür. Pik BPM genliği $\approx\SI{384}{\micro\meter}$'dir.

\item \textbf{BPM ofseti ham Fourier seviyesi ($\approx\SI{26}{\micro\meter}$).}
  $\mathbf{b}$ vektörü $R$'yi tamamen atlayarak BPM ölçümlerine doğrudan girer
  ve beyaz dağılımlıdır. Bunun FODO-antisimetrik $F_k$ moduna projeksiyonu
  $\sigma_b\sqrt{\pi/48}\approx\SI{26}{\micro\meter}$ r.m.s.\ değerinde olup
  tüm $k$'lara eşit dağılır; $k=2$'de özel bir tepe yoktur.
  Sinyal yörüngesi ($\SI{1669}{\micro\meter}$) bu ham Fourier seviyesinin
  zaten $64\times$ üzerindedir.

\item \textbf{Tahminleyici kirletme tabanı ($\approx\SI{0{,}6}{\micro\meter}$).}
  Tahminleyici $\mathbf{b}$'yi $F_{k=2}$'ye değil, $M_{k=2}=RF_{k=2}$'ye
  (yani $R$'nin $k=2$ sütununa) yansıtır ve $\|M_{k=2}\|=167$'ye böler:
  $\sigma_b/\|M_{k=2}\|=\SI{100}{\micro\meter}/167\approx\SI{0{,}6}{\micro\meter}$.
\end{enumerate}

Bu üç ölçek arasındaki hiyerarşi
\begin{equation}
  \SI{1669}{\micro\meter}\;\gg\;\SI{26}{\micro\meter}\;\gg\;\SI{0{,}6}{\micro\meter}
\end{equation}
yörünge kuvvetlenmesinin tahminleyici için çift yönlü çalıştığını göstermektedir:
gerçek sinyali BPM uzayında büyütürken ofseti $\|M_{k=2}\|$ ile bölerek baskılar.

Eşleşmiş-filtre SNR'si\footnote{Bu hesapta $\sigma_b=\SI{100}{\micro\meter}$
kullanılmıştır; bu değer BBA öncesi tipik BPM ofseti büyüklüğüdür.
Bölüm~\ref{sec:hatalar}'daki sistematik hata analizinde kabul edilebilir
seviye olarak $\sigma_b=\SI{50}{\micro\meter}$ (BBA sonrası) alınmıştır;
bu durumda SNR $\approx 33{,}3$'e çıkar.} $\sigma_b=\SI{100}{\micro\meter}$ için:
\begin{equation}
  \mathrm{SNR} = \frac{A\cdot\|M_{k=2}\|}{\sigma_b}
               = \frac{10\times167}{100} \approx 16{,}7.
\end{equation}

\begin{figure}[ht]
\centering
\fbox{\parbox{0.75\textwidth}{\centering\vspace{2cm}
\textbf{[ŞEKİL 3]} \\
Üç genlik ölçeği: sinyal yörüngesi 1669~μm, ham Fourier ofseti ~26~μm,
tahminleyici kirletme tabanı ~0.6~μm. \\
Sol panel: BPM okumaları ($k=2$ sinyal + $\sigma_b=100$~μm ofset).
Sağ panel: eşleşmiş-filtre projeksiyonu — sinyal ve ofset ayrımı.
\vspace{2cm}}}
\caption{BPM ofseti bastırımının üç ölçeği. Eşleşmiş-filtre yapısı,
$\mathbf{b}$ vektörünü $k=2$ yörünge desenine yansıtarak ofseti
$\|M_{k=2}\|=167$ kat bastırır. Sinyali $\mathbf{b}$'den ayırt eden
özellik büyüklük değil; $k=2$ ``parmak izi'' ile örtüşme derecesidir.}
\label{fig:amplitudes}
\end{figure}

\textbf{BPM ofsetine karşı dayanıklılık (beyazlık varsayımsız).}
Fourier-tabanlı formülasyon ``$\mathbf{b}$ ne kadar büyük?'' sorusunu
``$\mathbf{b}$'nin $k=2$ yörünge desenine projeksiyonu ne kadar büyük?'' sorusuna
dönüştürür. $M=RF$ uyum matrisi için $\mathbf{b}$'nin $k=2$ genlik tahminini
bozması:
\begin{equation}
  \delta a_{k=2} = \frac{\mathbf{b}\cdot\hat{m}_{k=2}}{\|M_{k=2}\|}
                 = \frac{\mathbf{b}\cdot\hat{m}_{k=2}}{167},
  \label{eq:bias}
\end{equation}
burada $\hat{m}_{k=2}=M_{k=2}/\|M_{k=2}\|$ birim vektördür. Denklem~(\ref{eq:bias})'e
$\mathbf{b}$'nin istatistiksel yapısına dair hiçbir varsayım girmez. En kötü durum
olan $\mathbf{b}=A\hat{m}_{k=2}$ (ofsetin tamamı $k=2$ yörünge deseniyle hizalı)
senaryosunda bile $A=\SI{100}{\micro\meter}$ için sapma yalnızca $\SI{0{,}60}{\micro\meter}$'dir;
güvenli limit $\|\mathbf{b}_\parallel\|<\SI{1669}{\micro\meter}$ ise gerçekçi
değerin $16{,}7\times$ üzerindedir.

%==============================================================
\section{CLEAN algoritması ile geri çatım}
\label{sec:clean}

\subsection{Algoritmanın çalışma prensibi ve neden CLEAN?}

Hizalama hatası spektrumunda hedef modlar ($k=1,2,3$; $10$--$30\,\mu$m) ile
kirletici modlar ($k=4\ldots10$; $100$--$300\,\mu$m) arasında $10$--$30\times$
genlik farkı vardır. Bu koşulda doğrudan en küçük kareler çözümü (lstsq) işe
yarar; fakat gürültülü ya da az belirlenmiş problemlerde büyük kirleticilerden
sızan hata küçük modları bozar. CLEAN, bu sorunu iteratif olarak çözer:
\textit{her adımda en baskın bileşeni artıktan önce ayıklar, ardından bir
sonrakine geçer.} Büyük kirleticiler ilk olarak temizlendiğinden küçük hedef
modlar kirletici gölgesinde kalmaz. Bu sıralama, örtük en küçük karelerde
olmayan bir "kirlilik izolasyonu" sağlar ve özellikle BPM gürültüsünün yüksek
olduğu senaryolarda lstsq'ya göre üstün performans verir.

CLEAN, temel olarak radyo astronomi görüntüleme alanından uyarlanan
iteratif bir bileşen soyutlama algoritmasıdır. Bu çalışmada aşağıdaki
uyarlanmış biçimiyle kullanılmaktadır:

\begin{enumerate}
  \item Aday $k$ değerleri kümesi belirlenir ($k=1,\ldots,10$).
  \item Her iterasyonda, mevcut artık yörünge $\mathbf{r}$ ile en yüksek
    korelasyona sahip olan $k$ bileşeni tespit edilir:
    $k^* = \arg\max_k \|M_k^T \mathbf{r}\| / \|M_k\|^2$.
  \item $k^*$ katkısının bir kısmı ($\gamma=0{,}1$ kazanç faktörü) artıktan
    çıkarılır ve birikimli katsayı listesine eklenir.
  \item Artık \SI{1}{\micro\meter} eşiğinin altına düşene veya maksimum
    iterasyon sayısına ulaşılana kadar tekrarlanır.
  \item Son katsayılardan her $k$ için genlik ve faz hesaplanır.
\end{enumerate}

Bu yaklaşımın temel avantajı, büyük kirletici bileşenleri ($k=4\ldots10$,
genellikle $100$--$300\,\mu$m) küçük hedef bileşenlerden ($k=1,2,3$;
$10$--$30\,\mu$m) önce ayrıştırmasıdır. Böylece küçük genlikli modlar,
büyük genliklilerin gölgesinde kaybolmamaktadır.

\subsection{CLEAN, tam baz ve kısmi baz kıyası}

Üç geri çatım stratejisi kıyaslanmıştır:
\begin{description}
  \item[(A) Kısmi baz ($k=1,2,3$):] Yalnızca üç Fourier modunu içeren
    baz ile en küçük kareler çözümü. Büyük $k=4\ldots10$ kirleticileri
    kısmi baza sızar; $k=2$'de $+1{,}34\%$ sistematik sapma oluşur.
  \item[(B) Tam baz ($k=1\ldots10$) lstsq:] Tüm modları kapsayan en küçük
    kareler çözümü. Kirletici sızıntısı yoktur, $k=2$ için $\approx0\%$ hata.
  \item[(C) CLEAN ($k=1\ldots10$ aday):] İteratif bileşen soyutlama.
    Kirletici sızıntısı yoktur; gürültülü koşullarda lstsq'ya göre üstündür
    çünkü büyük bileşenleri önce izole eder.
\end{description}

\begin{figure}[ht]
\centering
\fbox{\parbox{0.75\textwidth}{\centering\vspace{2cm}
\textbf{[ŞEKİL 4]} \textit{(Geri çatım kıyası)} \\
Üç yöntemle geri çatılan $k=1,2,3$ genlikleri. \\
Gerçek değerler: k=1: 30~μm, k=2: 10~μm, k=3: 25~μm. \\
Kirleticiler: k=4..10, her biri 100--300~μm.
\vspace{2cm}}}
\caption{Geri çatım kalitesi karşılaştırması: kısmi baz (A), tam baz lstsq (B),
CLEAN (C). Kısmi baz kirletici sızıntısı nedeniyle sapma gösterirken
tam baz ve CLEAN her üç modu da başarıyla kurtarır.}
\label{fig:reconstruction}
\end{figure}

\begin{figure}[ht]
\centering
\fbox{\parbox{0.75\textwidth}{\centering\vspace{2cm}
\textbf{[ŞEKİL 5]} \\
CLEAN iterasyon basamakları görseli. \\
Üst sol: başlangıç artık yörüngesi (tüm modlar mevcut). \\
Sonraki paneller: her iterasyonda en baskın bileşen (k=6, k=8, \ldots, k=2) \\
ayrıştırıldıkça artığın nasıl küçüldüğü. \\
Son panel: artık $<1\,\mu$m eşiğinde, $k=1,2,3$ tahminleri yerleşmiş.
\vspace{2cm}}}
\caption{CLEAN iterasyon adımları. Büyük kirletici bileşenler ($k=4\ldots10$)
ilk olarak ayrıştırılır; küçük hedef modlar ($k=1,2,3$) ancak kirleticiler
temizlendikten sonra net biçimde ortaya çıkar. Bu sıralama, lstsq'nun
yetersiz kaldığı gürültülü koşullarda CLEAN'i üstün kılar.}
\label{fig:clean_iterations}
\end{figure}

\subsection{Fiziksel yorum}

Alan kaynağı yalnızca kuadrupol hizalama hatalarıysa, $k$'ıncı Fourier
bileşenini tespit etmek tüm o $k$'ya karşılık gelen kuadrupol kaçıklıklarını
tespit etmeye eşdeğerdir. Rezonans $k=2$--$3$ civarında olduğundan, yüksek
$k$'ları kabaca, $k=1,2,3$'ü ise hassas biçimde bilmek yeterlidir: sahte
EDM'ye katkı $k\ge4$'ten giderek azalmaktadır (Şekil~\ref{fig:falseedm}).

%==============================================================
\section{Monte Carlo sistematik hata analizi}
\label{sec:hatalar}

Dört ayrı enstrümantal hata kaynağının $k=1,2,3$ modları üzerindeki etkisi
Monte Carlo simülasyonu ($N=200$ realizasyon) ile sayısal olarak incelenmiştir.
Gerçek hizalama deseni: $k=1$ ($30\,\mu$m), $k=2$ ($10\,\mu$m), $k=3$ ($25\,\mu$m)
artı $k=4\ldots10$ rastgele kirleticiler ($100$--$300\,\mu$m). Tüm geri çatımlar
CLEAN algoritması ile $k=1\ldots10$ aday setiyle yapılmıştır.

\subsection{BPM okuma gürültüsü}

Her ölçümde bağımsız Gaussian gürültü:
$\varepsilon_i\sim\mathcal{N}(0,\sigma^2)$.
$\sigma=5\,\mu$m'de $k=2$ genlik belirsizliği $\pm0{,}32\%$;
$\sigma=50\,\mu$m'de $\pm5{,}2\%$'ye ulaşır.
Analitik ölçekleme $\sigma(\Delta A/A)\propto\sigma/\|M_k\|$ simülasyonla
doğrulanmıştır.

\subsection{BPM statik ofseti}

Her BPM'e ölçümden ölçüme sabit, BPM'den BPM'e rastgele önyargı:
$b_i\sim\mathcal{N}(0,\sigma_b^2)$.
Analitik öngörü $\sigma(\Delta A/A)\approx\sigma_b/\|M_k\|$ simülasyonca
onaylanır: $\sigma_b=50\,\mu$m için $k=2$'de r.m.s.\ hata $0{,}4\%$,
$\sigma_b=200\,\mu$m için $1{,}5\%$.

\subsection{Kuadrupol rulosu}

$\theta$ açısıyla yuvarlanan bir kuadrupol, skew-quad terimi aracılığıyla
yatay kapalı yörüngeyi düşeye karıştırır:
$\delta\mathrm{KYB}\approx R_{dy}\cdot(2\theta\cdot x_\mathrm{co})$,
$x_\mathrm{co}=R_{dx}\cdot\delta x$, $\sigma_{dx}=100\,\mu$m.
$\theta_\mathrm{rms}=1\,\mathrm{mrad}$'da $k=2$ genlik hatası $-0{,}66\pm2{,}1\%$.

\subsection{Kuadrupol gradyan hataları}

Element başına fraksiyonel gradyan hataları $\varepsilon_j\sim\mathcal{N}(0,\sigma_G^2)$,
sütun ölçekleme modeliyle tepki matrisini bozar:
$\mathbf{y}_\mathrm{ölç}\approx\mathbf{y}_\mathrm{nom}+R\cdot(\varepsilon\odot\Delta y_\mathrm{gerçek})$.
$\sigma_G=0{,}5\%$'de $k=2$ hatası $-0{,}84\pm3{,}9\%$;
$\sigma_G=2\%$'de $\pm16\%$'ya çıkar, bu nedenle sıkı gradyan kalibrasyonu önem taşır.

\subsection{Kombinat hata bütçesi}

Dört etki kabul edilebilir seviyelerinde ($\sigma_\mathrm{gürültü}=5\,\mu$m,
$\sigma_b=50\,\mu$m, $\theta_\mathrm{rms}=1\,\mathrm{mrad}$,
$\sigma_G=0{,}5\%$) eş zamanlı uygulandığında, Tablo~\ref{tab:hatalar}'da
özetlenen sonuçlar elde edilir.

\begin{table}[h]
\centering
\caption{Kabul edilebilir seviyelerde sistematik hata bütçesi (CLEAN, $N=200$ MC).
$\mu$ ve $\sigma$: Monte Carlo topluluğu üzerinden $\Delta A/A$ [\%]'nin
ortalama ve standart sapması.}
\label{tab:hatalar}
\begin{tabular}{lcccc}
\toprule
Kaynak & Seviye & $k=1$ ($\mu\pm\sigma$) & $k=2$ ($\mu\pm\sigma$) & $k=3$ ($\mu\pm\sigma$) \\
\midrule
BPM gürültüsü    & $5\,\mu$m            & $-0{,}93\pm0{,}38$ & $-0{,}74\pm0{,}32$ & $-0{,}92\pm0{,}49$ \\
BPM ofseti       & $50\,\mu$m           & $-0{,}80\pm3{,}98$ & $-0{,}13\pm2{,}51$ & $-0{,}66\pm4{,}69$ \\
Kuadrupol rulosu & $1\,\mathrm{mrad}$   & $-0{,}89\pm0{,}62$ & $-0{,}62\pm1{,}92$ & $-0{,}95\pm0{,}72$ \\
Gradyan hatası   & $0{,}5\%$            & $-0{,}90\pm1{,}45$ & $-0{,}84\pm3{,}94$ & $-1{,}15\pm1{,}71$ \\
\midrule
Kombinat         & tümü                 & $-0{,}83\pm4{,}38$ & $-0{,}12\pm5{,}37$ & $-0{,}69\pm5{,}21$ \\
\bottomrule
\end{tabular}
\caption*{$k=1,2,3$ için kombinat belirsizlik $\pm4$--$5\%$ düzeyinde;
$\lesssim10\%$ göreli genlik doğruluğu hedefi karşılanmaktadır.
Baskın kaynak: BPM ofseti ($k=1,3$ için) ve gradyan hatası ($k=2$ için).
$\sigma_G\lesssim0{,}2\%$ kalibrasyonuyla kombinat belirsizlik $\sim2\%$'ye iner.}
\end{table}

\begin{figure}[ht]
\centering
\fbox{\parbox{0.75\textwidth}{\centering\vspace{2cm}
\textbf{[ŞEKİL 6]} \textit{(Kombinat hata bütçesi)} \\
$k=1,2,3$ için kombinat sistematik hata çubuk grafiği. \\
Her panel bir mod; yatay eksen: hata kaynağı (BPM gürültüsü, ofset, rulo, gradyan, kombinat). \\
Dikey eksen: $\Delta A/A$ [\%] (ortalama ± 1σ). N=200 MC.
\vspace{2cm}}}
\caption{Sistematik hata katkıları, $k=1,2,3$ modları için ayrı ayrı.
Her çubuk: o hata kaynağının tek başına oluşturduğu etki; son çubuk kombinattır.}
\label{fig:sistematics}
\end{figure}

%==============================================================
\section{Harici radyal alan harmoniklerine genişleme}
\label{sec:haricalan}

\subsection{Yöntemin kapsamı}

Bu yöntem, kapalı-yörünge bozunumunun $k$'ıncı harmonikini \textit{kökeninden
bağımsız} olarak ölçer. Halka boyunca yayılan harici bir radyal manyetik alan
$B_x(\theta)=B_r\cos(k\theta)$, $k$'ıncı harmonik kuadrupol kaçıklığıyla tamamen
aynı Fourier bileşeninde yörünge bozunumu yaratır; çünkü ikisi de spini aynı
mekanizma üzerinden etkiler: halka genelinde toplanan radyal impuls.
Dolayısıyla yöntem, mekanik hizalama hatalarını ve harici manyetik alan
harmoniklerini (sistematik alan kusurları, artık saçak alanlar vb.) aynı
çerçevede kapsayarak ölçer.

\subsection{Spin-eşdeğer radyal alan ve dönüşüm faktörü}

Duyarlılık sonuçlarını fiziksel alan birimi cinsinden ifade etmek için
spin-eşdeğer radyal alan $B_{x,\mathrm{eq}}^{(k)}$ tanımlanır:
$A_q$ genlikli $k$'ıncı harmonik kaçıklıkla aynı $dS_y/dt$'yi üretecek
$B_x(\theta)=B_{x,\mathrm{eq}}\cos(k\theta)$ genliği.
Dönüşüm faktörü:
\begin{equation}
  c_k \approx 6\,\mathrm{nT}\,\mu\mathrm{m}^{-1}
  \quad (k=1,\ldots,7,\;\pm17\%)
  \label{eq:ck}
\end{equation}
neredeyse $k$'dan bağımsızdır; spin, lokal pik alan yerine halka genelindeki
ortalama impulse tepki verdiğinden FODO dolgu faktörü $f_\mathrm{fill}\sim3\%$
ve gradyan $G=\SI{0{,}21}{\tesla\per\meter}$ genel ölçeği belirler.
Lokal kuadrupol-merkezi alanı $B_{x,\mathrm{eq}}$'nun $\sim\!34\times$'idir.

\textbf{Dayanıklılık notu.} $c_k$'nın $k$'dan bağımsızlığı pratik açıdan
kritik bir özelliktir: hangi mod ölçülürse ölçülsün ($k=1,2$ veya $3$),
aynı $\approx6\,\mathrm{nT}\,\mu\mathrm{m}^{-1}$ dönüşüm katsayısı
kullanılabilir. Bu, modele özgü kalibrasyona gerek olmadığı ve çerçevenin
kafes değişikliklerine ve farklı deney koşullarına karşı güçlü olduğu
anlamına gelir.

\subsection{Ölçüm hassasiyeti}

BPM sistematik ofsetinden kaynaklanan alan belirsizliği:
\begin{equation}
  \sigma(B_{x,\mathrm{eq}}^{(k)}) \approx
  \frac{c_k\,\sigma_b}{\|M_k\|}.
\end{equation}

1~nT hassasiyet için gereken BPM ofseti:
\begin{itemize}
  \item $k=2$: $\sigma_b\lesssim\SI{21}{\micro\meter}$ (BBA sonrası ulaşılabilir)
  \item $k=3$: $\sigma_b\lesssim\SI{5}{\micro\meter}$ (BBA hedefi)
\end{itemize}

Kuadrupoller ve diğer manyetik alan kaynakları $\sim\SI{20}{\micro\meter}$
düzeyine hizalandığında (BBA sonrası), $k=2$ bileşeni $\sim\SI{1}{\nano\tesla}$,
$k=3$ bileşeni ise $\sim\SI{4}{\nano\tesla}$ hassasiyetle ölçülebilmektedir.

%==============================================================
\section{Tartışma ve sonraki çalışmalar}
\label{sec:tartisma}

Sunulan yöntem, k-modülasyonunun ofset-sıfırlama garantisi ile tek-yörünge
basitliğini özellikle baskın sahte-EDM harmonik için etkin olan yapısal
bir ayrıştırmayla takas eder. Çalışmanın kapsamı kasıtlı olarak hedeflidir:
tam hizalama vektörünü geri çatmak değil, tek fiziksel gözlenebiliri
($k=2$ genliği ile birlikte $k=1$ ve $k=3$) ölçmek.

\textbf{Yöntemin operasyonel sınırı.}
Bu yöntem $k=1,2,3$ harmonik genliklerini \emph{ölçer}; bunların düzeltici
mıknatıslarla bastırılmasını ise kapsamaz. Pratikte ölçüm-düzeltme döngüsü
tekrarlanacaktır: $k=2$ düzeltildikten sonra biriken artık harmonik, bir
sonraki orbit ölçümüyle izlenebilir. Yeterince küçük artık için aynı
çerçeve geçerliliğini korur; çünkü eşleşmiş-filtre SNR'si artık büyüklüğüyle
orantılı olarak küçülse de BPM-ofset bastırımı mekanizması değişmez.

Aşağıdaki konular sonraki çalışmalar için açık kalmaktadır:

\begin{itemize}
  \item \textbf{Gerçek makine verileriyle doğrulama.}
    Bu çalışmada Omarov et al.~\cite{omarov} pEDM halka tasarım kafesi
    ($N_\mathrm{FODO}=24$, $Q_y=2{,}68$, $G=\SI{0{,}21}{\tesla\per\meter}$,
    $R_0=\SI{95{,}49}{\meter}$) kullanılmıştır; dolayısıyla ``idealize FODO''
    itirazı büyük ölçüde geçersizdir.
    Açık kalan doğrulama adımı, bir prototip ya da gerçek pEDM depolama
    halkasından alınan BPM verileriyle yöntemin deneysel olarak sınanmasıdır.
  \item \textbf{Çoklu yörünge ortalaması.}
    Farklı zamanlarda alınan yörüngelerin birleştirilmesi, yavaş değişen
    $\mathbf{b}$'yi azaltırken $\Delta q$'yu sabit tutar; ofseti daha da bastırır.
  \item \textbf{Tepki matrisi kalibrasyonu.}
    $R$'nin LOCO/beta-BBA ile ne kadar iyi ölçülebildiği ve kalan
    $\delta K/K$'nın $k=2$ hata bütçesine yayılması incelenmelidir.
  \item \textbf{Gradyan kalibrasyonu — en yüksek öncelikli donanımsal iyileştirme.}
    Monte Carlo analizinde $\sigma_G=0{,}5\%$, üç mod için de baskın belirsizlik
    kaynağıdır ($k=2$'de $\pm3{,}9\%$, kombinat $\pm5{,}4\%$). Buna karşın BPM
    ofseti ve kuadrupol rulosu yeterince küçük tutulsa bile tek başına gradyan
    hatası bu sınırlamayı koruyacaktır. LOCO/beta-BBA tabanlı in-situ gradyan
    kalibrasyonu ile $\sigma_G\lesssim0{,}2\%$'ye ulaşılırsa kombinat belirsizlik
    $\sim2\%$'ye iner ve alan hassasiyeti de buna bağlı olarak $\sim3\times$
    iyileşir. Bu, deney grubunun en yüksek getiri sağlayacağı tek kalibrasyon
    adımıdır.
  \item \textbf{Yatay düzlem ve kuadrupol tilti.}
    Halihazırda oluşturulan yatay tepki matrisi $R_x$ kullanılarak yatay
    düzleme ve skew kuplaja genişleme yapılabilir.
  \item \textbf{Harmoniğin spin düzeyine bağlanması.}
    Ölçülen $\Delta q_{k=2}$'nin yanlış-EDM fonksiyoneline
    $\Phi_\mathrm{spin}\propto\sum_j K_jL_j(y_\mathrm{KY}(s_j)-\Delta y_j)$
    oranlanması ve kapatma toleransının tanımlanması gerekmektedir.
\end{itemize}

%==============================================================
\section{Sonuç}
\label{sec:sonuc}

Bu çalışmada, pEDM dondurulmuş-spin halkasında sahte EDM'ye yol açan baskın
hizalama hatası harmoniklerinin ($k=1,2,3$) tek bir kapalı-yörünge ölçümünden
CLEAN algoritması kullanılarak $\sim5\%$ göreli genlik doğruluğuyla
kurtarılabildiği simülasyonla gösterilmiştir.

Yöntemin temel dayanağı, tepki matrisinin üç özelliğidir:
\begin{itemize}
  \item $k=2$ için rezonans kuvvetlenmesi: $\|M_{k=2}\|=167$, kafes
    ortalamasının $\sim34$ katı; bu sayede $\SI{10}{\micro\meter}$'lik
    hizalama hatası BPM uzayında $\SI{1669}{\micro\meter}$ yörünge normuna dönüşür.
  \item Eşleşmiş-filtre yapısı, $\SI{100}{\micro\meter}$'lik BPM ofsetini
    $\SI{0{,}6}{\micro\meter}$ kirletme tabanına bastırır —
    $\mathbf{b}$'nin istatistiksel yapısına dair herhangi bir varsayım olmaksızın.
  \item Farklı Fourier modlarının yörünge tepkilerinin neredeyse dikeyliği,
    büyük kirletici harmoniklerin ($k=4\ldots10$, $100$--$300\,\mu$m) küçük
    hedef modlarına ($k=1,2,3$, $10$--$30\,\mu$m) sızmasını engeller.
\end{itemize}

Monte Carlo sistematik analizi, makul inşaat hataları koşullarında
($\sigma_b\lesssim50\,\mu$m, $\theta_\mathrm{rms}\lesssim1\,\mathrm{mrad}$,
$\sigma_G\lesssim0{,}5\%$) $\lesssim10\%$ göreli genlik doğruluğu hedefinin
üç mod için de karşılandığını doğrulamıştır. BBA sonrası
$\sigma_b\lesssim\SI{20}{\micro\meter}$ koşulunda $k=2$ moduna karşılık gelen
radyal alan harmonik bileşeni $\sim\SI{1}{\nano\tesla}$ hassasiyetle
ölçülebilmektedir.

%==============================================================
\section*{Teşekkür}
\todo{teşekkür / finansman}

\begin{thebibliography}{9}
\bibitem{omarov} Z.~Omarov et al., \emph{Comprehensive symmetric-hybrid
ring design for a proton EDM experiment at below $10^{-29}\,e\cdot$cm},
Phys.\ Rev.\ D \textbf{105}, 032001 (2022).
\bibitem{safranek1997} J.~Safranek, \emph{Experimental determination of storage
ring optics using orbit response measurements},
Nucl.\ Instrum.\ Methods Phys.\ Res.\ A \textbf{388}, 27 (1997).
\bibitem{todo-pEDM} \todo{pEDM deney önerisi referansı}
\end{thebibliography}

\end{document}
